\section{Chow niveau, complete intersections, and flat fibrations}\label{sec:niveau}

\subsection{A residual motive of small dimension}

We use covariant rational Chow motives, with
$\CH_a(h(T)(u))=\CH_{a-u}(T)_\Q$.

\begin{theorem}[Chow-niveau transfer]\label{thm:niveau}
Let $X/k$ have dimension $n$. Suppose that after a finite extension
there is an actual rational Chow decomposition
\begin{equation}\label{eq:niveau-decomposition}
 h(X)\simeq Q\oplus(T,e,u),\qquad
 Q\text{ of abelian type},\quad \dim T\leq n-2u,
\end{equation}
where $u\geq0$. Over a finite extension of $\Q_p$ this gives zero
obstruction on the full geometric homological Chow groups for
\begin{equation}\label{eq:niveau-local}
 r\leq u+1\quad\text{or}\quad r\geq n-u-1.
\end{equation}
Every degree is covered if $n\leq2u+3$. Over a number field the full
height pairs are obtained in every degree in that case; otherwise
they are obtained for
\begin{equation}\label{eq:niveau-height}
 r\leq u+1\quad\text{or}\quad r\geq n-u.
\end{equation}

A sufficient hypothesis for \eqref{eq:niveau-decomposition} is a
morphism $N\to h(X)$ from an abelian-type Chow motive which, on an
algebraically closed universal domain $\Omega/k$ of infinite
transcendence degree over the field of definition, surjects on
algebraically trivial $\CH_i$ for every $i<u$.
\end{theorem}

\begin{proof}
Let $m=\dim T$. The residual image of a codimension-$r$ class has
dimension $n-r-u$ and codimension
\begin{equation}\label{eq:residual-degree}
 m-n+r+u\leq r-u.
\end{equation}
If $r\leq u+1$, this is a divisor degree or lower. If
$r\geq n-u-1$, the residual cycle dimension is at most one.
The endpoint input applies in either case. If $n\leq2u+3$, then
$m\leq3$, so every residual degree is covered. The abelian-type
summands are covered by \thmref{thm:main}. Their actual inclusion
and retraction maps, followed by \lemref{lem:general-transfer},
prove the local assertion. Intersect its range with the range of
$n+1-r$ to obtain \eqref{eq:niveau-height}.

For the sufficient hypothesis apply Vial's
\cite[Theorem 3.11]{Vial} with algebraic equivalence and $M=h(X)$.
The finite-dimensionality required there holds for $N$. Self-duality
of $h(X)$ gives \eqref{eq:niveau-decomposition} over the field of
definition of the morphism. Its additional curve and Tate summands
are of abelian type. A negative residual dimension means that the
residual summand is zero.
\end{proof}

In particular it suffices that
$\CH_i(X_\Omega)_{\Q,\alg}=0$ for $0\leq i<u$: take $N=0$.
This criterion concerns the subgroup algebraically equivalent to
zero, rather than a quotient by algebraic equivalence.

\subsection{Small-degree Chow bounds}

Let $X\subset\PP^N$ be a smooth complete intersection of degrees
$d_1\geq\cdots\geq d_c\geq2$, and put $n=N-c$.

\begin{proposition}\label{prop:small-degree}
Let $0\leq t\leq n-1$. Suppose either
\begin{equation}\label{eq:elv}
 \left(d_1\geq3\ \text{or}\ c\geq t+1\right),\qquad
 \sum_{a=1}^c\binom{d_a+t}{t+1}\leq N,
\end{equation}
or all $d_a=2$, $1\leq c\leq t$, and
\begin{equation}\label{eq:elv-quadratic}
 c(t+2)\leq N+c-t-1.
\end{equation}
Then \thmref{thm:niveau} applies with $u=t+1$.
For a hypersurface of degree $d\geq3$, an alternative sufficient
condition is
\begin{equation}\label{eq:hi-bound}
 (j+1)(N-j)+1\geq\binom{d+j+1}{j+1}
             \quad(0\leq j\leq t).
\end{equation}
In each case all degrees are covered if $n\leq2t+5$; otherwise
the guaranteed full height degrees are $r\leq t+2$ or
$r\geq n-t-1$.
\end{proposition}

\begin{proof}
Theorem 4.6 of Esnault--Levine--Viehweg \cite{ELV} gives
$\CH_i(X_\Omega)_\Q=\Q$ for $0\leq i\leq t$ under
\eqref{eq:elv} or \eqref{eq:elv-quadratic}. In the hypersurface
case, Hirschowitz--Iyer \cite[Theorem 1.5]{HI} gives the same
conclusion degree by degree under \eqref{eq:hi-bound}; the upper
variety is projective space and its relevant Chow group is
one-dimensional. A linear-section generator has nonzero projective
degree. The algebraically trivial subgroup is therefore zero in
all these degrees. Apply \thmref{thm:niveau} with $u=t+1$.
\end{proof}

There is a further useful bound for mixed quadrics and cubics.
Otwinowska's rewritten author version \cite[Corollaire 1]{Otwinowska}
gives the following statement for smooth complete intersections:
if there are $a$ quadrics and $b\geq1$ cubics in $\PP^N$, $N\leq15$,
then $\CH_i(-)_\Q$ is generated by the ambient section class for
\begin{equation}\label{eq:ot-bound}
 0\leq i\leq t=\left\lfloor\frac{N-2a-3b}{3}\right\rfloor.
\end{equation}
Its complex-field formulation gives the required universal-domain
statement in characteristic zero. Indeed, descend a given cycle
and its equations to a finitely generated field, embed that field
in $\C$, and obtain the rational-equivalence certificate there.
The certificate descends to a further finitely generated field,
which embeds over the original field of definition into the
universal domain. Its base change proves the original relation.
The coefficient of the ambient class is fixed by projective degree.

\begin{corollary}[Explicit complete-intersection domains]\label{cor:ci-list}
Every smooth geometrically integral variety in the following table
has $\mathsf V$ locally and a full height domain over a number
field. A presentation after a finite extension suffices.
\begin{center}
\begin{tabular}{lll}
\toprule
Type & Dimension & Low-Chow parameter $t$\\
\midrule
Cubic hypersurface & $6,7$ & $1$\\
Cubic hypersurface & $8,9$ & $2$\\
Cubic hypersurface & $11$ & $3$\\
Quartic hypersurface & $7$ & $1$\\
Complete intersection $(2,3)$ & $6,7$ & $1$\\
Complete intersection $(2,3)$ & $9$ & $2$\\
Complete intersection $(2,2,3)$ & $7$ & $1$\\
Complete intersection $(3,3)$ & $7$ & $1$\\
\bottomrule
\end{tabular}
\end{center}
The cubic cases in dimension at most five already follow from
\corref{cor:rc-fibrations}.
\end{corollary}

\begin{proof}
For cubics the values in \eqref{eq:hi-bound} are
$4,10,20,35$ for $j=0,1,2,3$. The requisite minimum ambient
dimensions for $t=1,2,3$ are $6,9,12$, respectively. These imply
the first three rows. For the quartic row, $N=8$ and the two
inequalities are $9\geq5$ and $15\geq15$.

For $(2,3)$ of dimension $6$ or $7$, the zero-cycle statement
follows from \eqref{eq:elv}. The upper quadric, which may be
singular, has $\CH_2=\Q$ by the all-intersection form of
\cite[Theorem 4.6]{ELV} and \eqref{eq:elv-quadratic}.
The distinct-largest-degree
case of \cite[Theorem 1.5]{HI}, with $j=1,c=2$, has
$\rho=3N-19$ and $\rho+1\geq N-2$ exactly when $N\geq8$.
Thus $\CH_1=\Q$. The last three rows follow from
\eqref{eq:ot-bound}: their ambient dimensions are $11,10,9$,
respectively. In every row $n\leq2t+5$.
\end{proof}

\subsection{The all-fiber criterion}

\begin{proposition}[Flat-fibration transfer]\label{prop:fibration}
Let $f:Y\to B$ be flat, projective and surjective, with $Y,B$
smooth projective and $\dim Y=n$, $\dim B=b$. Fix an integer
$u\geq1$ with $n-2u\geq0$. Suppose that for
every geometric scheme fiber over a universal domain
\[
 \CH_i(Y_\omega)_\Q=\Q\quad(0\leq i<u).
\]
Assume $\mathsf V(B)$ in the native local theory. Then the full
local source vanishes for $r\leq u+1$ or $r\geq n-u-1$, and in
all degrees if $n-2u\leq3$. If these base hypotheses hold at every
finite place of a number-field model, the corresponding full
height degrees are obtained by \eqref{eq:range-intersection}.
\end{proposition}

\begin{proof}
Vial's flat-fibration decomposition \cite[Theorem 4.2]{VialFibrations}
is an actual Chow decomposition
\begin{equation}\label{eq:fibration-motive}
 h(Y)\simeq\bigoplus_{j=0}^{n-b}h(B)(j)\oplus(T,e,u),
 \qquad\dim T=n-2u.
\end{equation}
The base summands are admitted by hypothesis. The residual-degree
calculation \eqref{eq:residual-degree} and the endpoint input give
the stated local degrees. Transfer through the actual diagonal
and apply \lemref{lem:two-degrees} globally.
\end{proof}

\begin{corollary}[Quadric and cubic fibers]\label{cor:quad-cubic-fibers}
The following flat projective surjective maps, with smooth projective
geometrically integral total space and base, have $\mathsf V$:
\begin{enumerate}[label=\textup{(\roman*)}]
\item Every geometric scheme fiber is a nonzero quadric of fixed
dimension $q\geq1$ in $\PP^{q+1}$, and $1\leq b\leq3$.
The same holds for $b=4$, $q$ odd, if the base has $\mathsf V$;
in particular the base may be $\PP^4$.
\item Every geometric scheme fiber is a nonzero cubic hypersurface
of fixed dimension $q$, and
\[
 (q,b)=(5,1),(5,2),(6,1),\text{ or }(8,1).
\]
\end{enumerate}
The conclusions concern the entire geometric homological Chow
groups. Over number fields the same geometric hypotheses give
full height domains, with the base condition imposed at all
finite places when $b=4$.
\end{corollary}

\begin{proof}
A nonzero quadric scheme of dimension $q$ over an algebraically
closed characteristic-zero field has $\CH_i\simeq\Z$ for
$i<q/2$. This includes rank-one quadrics: Chow groups are
unchanged by reduction. For other singular quadrics, the cone
description and the split-quadric cellular filtration give the
same range; the cone localization proof is recorded in
\cite{InternalFibrations}. Take $u=\lceil q/2\rceil$ in
\propref{prop:fibration}. Its residual dimension is $b$ for even
$q$ and $b-1$ for odd $q$, giving (i).

For cubics of dimension at least five, the all-intersection
ELV bound gives $\CH_0=\CH_1=\Q$. For $q=8$, the HI bound
also gives $\CH_2=\Q$. These cited bounds apply to the reduced
singular hypersurfaces in question. A nonreduced cubic has
reduction a hyperplane or a union of two hyperplanes; the closed
union Chow sequence gives the same low-degree groups. Thus the
four cases have residual dimensions $2,3,3,3$, respectively,
using $u=2,2,2,3$. The base has dimension at most two and hence
has $\mathsf V$. Apply \propref{prop:fibration}.
\end{proof}

\subsection{A signed retraction for intersections of quadrics}

\begin{theorem}[Intersections of quadrics]\label{thm:quadrics}
A smooth complete intersection of $c$ quadrics in $\PP^N$,
$1\leq c\leq N-1$, has $\mathsf V$ in all degrees when
\[
 c\leq4,\qquad\text{or}\qquad c=5\text{ and }N\text{ is even}.
\]
For arbitrary $c$, the local degrees supplied are
\begin{equation}\label{eq:quadrics-local}
 r\leq\lfloor N/2\rfloor-c+2
 \quad\text{or}\quad r\geq\lceil N/2\rceil-2.
\end{equation}
The full global height degrees outside the full case are
$r\leq\lfloor N/2\rfloor-c+2$ or
$r\geq\lceil N/2\rceil-1$.
\end{theorem}

\begin{proof}
For $c=1$, pass to a finite splitting field of the smooth quadric.
Its cellular Chow decomposition is Tate, so \thmref{thm:main}
gives every degree. Assume $c\geq2$. Write the equations as
$Q_1,\ldots,Q_c$ and let
\[
 Y=\{\textstyle\sum_jt_jQ_j=0\}\subset\PP^N\times\PP^{c-1}.
\]
It is a flat family of nonzero $(N-1)$-dimensional quadric
schemes. The total space is smooth: at a possible singular point
all $Q_j$ vanish, while their independent conormal differentials
on the smooth complete intersection make $\sum_jt_jdQ_j$
nonzero. Set $m=c-1$, $E=X\times\PP^m$, and write
$j:E\hookrightarrow Y$, $p:E\to X$. The normal sequence is
\[
 0\longrightarrow N_{E/Y}\longrightarrow
 p^*\mathcal O_X(2)^{\oplus c}\longrightarrow
 p^*\mathcal O_X(2)\otimes\mathcal O_{\PP^m}(1)
 \longrightarrow0.
\]
Its degree-$m$ Chern class has fiber-leading term
$(-1)^mc_1(\mathcal O_{\PP^m}(1))^m$, and fiber integration gives
\[
 p_*c_m(N_{E/Y})=(-1)^m.
\]
Consequently the correspondences
\begin{equation}\label{eq:quadrics-retraction}
 U=j_*p^*,\qquad V=(-1)^mp_*j^!,\qquad VU=1
\end{equation}
give a retraction, by self-intersection. The calculation after
every smooth projective product shows the same identity for
Chow correspondences. Here $U$ sends codimension $r$ to $r+m$.

Apply \propref{prop:fibration} to $Y\to\PP^{c-1}$, with
$u=\lceil(N-1)/2\rceil$. Its residual dimension is $c-2$ if
$N$ is even, and $c-1$ if $N$ is odd. These are at most three
in the asserted full cases. Substituting $r+m$ in the outer
range and using \eqref{eq:quadrics-retraction} gives
\eqref{eq:quadrics-local}. Intersecting the two complementary
ranges gives the height statement.
\end{proof}

\begin{theorem}[Relative quadrics]\label{thm:relative-quadrics}
Let $B$ be smooth projective geometrically integral of dimension
$b\geq1$, and let $\mathcal V,\mathcal F$ be vector bundles of
ranks $N+1,c$, with $c\geq1$ and $N\geq c+1$.
Let $\pi:\PP_B(\mathcal V)\to B$ parametrize lines. Suppose that
the zero scheme of a regular section of
$\mathcal O(2)\otimes\pi^*\mathcal F$ is smooth projective
geometrically integral, is flat over $B$, and has expected
codimension $c$ in every geometric fiber. If $\mathsf V(B)$,
then it has $\mathsf V$ whenever
\[
 b+c\leq5\quad(N\text{ even}),\qquad
 b+c\leq4\quad(N\text{ odd}).
\]
In general the local degree range is
\[
 r\leq\lfloor N/2\rfloor-c+2
 \quad\text{or}\quad r\geq b+\lceil N/2\rceil-2.
\]
Over a number field with the base admitted at every finite place,
the corresponding full height range ends instead at
$r\geq b+\lceil N/2\rceil-1$. The base condition is automatic
for $1\leq b\leq3$.
\end{theorem}

\begin{proof}
Form $C=\PP_B(\mathcal F^\vee)$ and the universal quadric
$Y\to C$. Expected fiber codimension ensures that no nonzero
coefficient vector gives the zero quadratic equation. Smoothness
of the total space follows from the absolute differential of the
regular section defining $X$; the base directions are included.
The projective-bundle formula gives $\mathsf V(C)$. For
$E=X\times_BC$, the same normal sequence and fiber integration
as in \eqref{eq:quadrics-retraction} give a signed retraction with
shift $c-1$. The residual dimension of the universal quadric is
$b+c-2$ for even $N$ and $b+c-1$ for odd $N$. Apply the
flat-fibration theorem, transfer through the retraction, and use
\lemref{lem:two-degrees}.
\end{proof}
